On considère la fonction $f$ définie pour tout réel $x$ par $f(x) = -x^2 + 6x - 5$.

\begin{enumerate}
\item Calculer l'image de 0 et de 3 par la fonction $f$.

\item Montrer que, pour tout réel $x$, on a : $(x - 1)(5 - x) = -x^2 + 6x - 5$.

\item En déduire les antécédents de 0 par la fonction $f$.

\item Montrer que pour tout réel $x$, on a : $4 - (x - 3)^2 = -x^2 + 6x - 5$.

\item Est-il possible de trouver un réel $x$, tel que $f(x) > 4$ ? Justifier.

\item Réaliser un schéma donnant l'allure de la courbe de la fonction $f$ sur lequel apparaîtront les résultats des questions \textbf{1.}, \textbf{3.} et \textbf{5.}
\end{enumerate}
